% Diagrams for the template manuscript. Each tikzpicture becomes one page (standalone, multi); make_figures.py turns the
% pages into fig_d1.png ... fig_d7.png.
\documentclass[multi=tikzpicture,border=5pt]{standalone}
\usepackage{amsmath}
\usepackage{mathptmx}
\usepackage[T1]{fontenc}
\usepackage{amssymb}
\usepackage{tikz}
\usetikzlibrary{arrows.meta,positioning,calc,shapes.geometric,fit,backgrounds,decorations.pathreplacing}
\definecolor{cIdx}{HTML}{A9B0E4}\definecolor{cRec}{HTML}{4A54C1}\definecolor{cUndo}{HTML}{E2B77A}
\definecolor{cFree}{HTML}{ECEAE4}\definecolor{cOk}{HTML}{CFE6D8}\definecolor{cBad}{HTML}{F3C9C1}
\definecolor{cAcc}{HTML}{E6E8F8}
\begin{document}
\tikzset{
  box/.style={draw,rounded corners=2pt,thick,align=center,font=\small,inner sep=4pt,fill=white},
  tb/.style={draw,thick,align=center,font=\footnotesize,inner sep=3pt,fill=white,rounded corners=1.5pt},
  arr/.style={-{Stealth[length=2.6mm]},thick},
  cell/.style={draw,minimum size=7.2mm,inner sep=0pt,font=\footnotesize,thick},
  lab/.style={font=\scriptsize,inner sep=1pt}
}

%%%%%%%%%%%%%%%%%%%%%%%%%%%%%%%%%%%%%%%% D1: architecture
\begin{tikzpicture}[x=1cm,y=1cm]
  \node[box,fill=cAcc,minimum width=3.6cm,minimum height=3.1cm,text width=3.3cm] (src) at (0,0)
    {\textbf{Compiler front end}\\[2pt]\footnotesize scope events from the source text\\[3pt]\texttt{declare(name,type)}\\\texttt{lookup(name)}\\\texttt{enter\_scope()}\\\texttt{exit\_scope()}};
  % buffer
  \node[draw,thick,rounded corners=2pt,minimum width=8.6cm,minimum height=3.1cm,fill=white] (buf) at (7.0,0) {};
  \node[font=\small\bfseries,anchor=south] at (7.0,1.62) {STRATA table in ONE fixed buffer (caller-provided)};
  \fill[cIdx] (2.95,0.45) rectangle (5.05,1.25);\draw[thick] (2.95,0.45) rectangle (5.05,1.25);
  \fill[cRec] (5.05,0.45) rectangle (7.5,1.25);\draw[thick] (5.05,0.45) rectangle (7.5,1.25);
  \fill[cFree] (7.5,0.45) rectangle (9.0,1.25);\draw[thick] (7.5,0.45) rectangle (9.0,1.25);
  \fill[cUndo] (9.0,0.45) rectangle (10.95,1.25);\draw[thick] (9.0,0.45) rectangle (10.95,1.25);
  \node[font=\scriptsize,align=center] at (4.0,0.85) {index slots\\open addressing};
  \node[font=\scriptsize,align=center,text=white] at (6.28,0.85) {identifier records\\(packed, one per name)};
  \node[font=\scriptsize] at (8.25,0.85) {free};
  \node[font=\scriptsize,align=center] at (9.98,0.85) {undo stack\\$\leftarrow$ grows};
  \node[font=\scriptsize,align=left,text width=7.9cm] at (7.0,-0.55)
    {Linear probing with probe limit $L$; a key never moves after it is placed.\\Space is reserved before anything is written: a request that does not fit changes nothing.};
  \node[box,fill=cAcc,minimum width=3.2cm,minimum height=3.1cm,text width=2.9cm] (out) at (14.0,0) {\textbf{Results}\\[2pt]\footnotesize status codes\\\texttt{kOk}, \texttt{kNoMemory},\\\texttt{kProbeLimit}, ...\\state digest};
  \draw[arr] (src.east) -- node[above,lab]{calls} (buf.west);
  \draw[arr] (buf.east) -- node[above,lab]{answers} (out.west);
  % analysis layer
  \node[font=\small\bfseries,anchor=west] at (-1.8,-2.55) {Analysis layer (this work)};
  \node[tb,minimum width=3.6cm,minimum height=1.25cm,text width=3.4cm] (a1) at (0.0,-4.0) {\textbf{Exact undo}\\Index restored exactly\\Safe deletion orders\\Orphan law};
  \node[tb,minimum width=3.6cm,minimum height=1.25cm,text width=3.4cm] (a2) at (4.1,-4.0) {\textbf{Refusal laws}\\Exact law (one push)\\Whole-fill law\\Unlimited run};
  \node[tb,minimum width=3.6cm,minimum height=1.25cm,text width=3.4cm] (a3) at (8.2,-4.0) {\textbf{Bounded churn}\\Persistent names,\\schedule-free\\bound $D_s$};
  \node[tb,minimum width=3.6cm,minimum height=1.25cm,text width=3.4cm] (a4) at (12.3,-4.0) {\textbf{Keys and space}\\Key lifetime and\\adversaries;\\space ledger};
  \node[box,fill=cOk,minimum width=15.6cm] (des) at (6.15,-6.15) {Design rules: probe limit $L$, buffer size, key policy};
  \foreach \a in {a1,a2,a3,a4} \draw[arr] (\a.south) -- (\a.south |- des.north);
  \draw[arr,dashed] (a2.north) -- (buf.south -| a2.north) node[midway,right,lab,align=left]{models the index};
\end{tikzpicture}

%%%%%%%%%%%%%%%%%%%%%%%%%%%%%%%%%%%%%%%% D2: exact undo
\begin{tikzpicture}[x=1cm,y=1cm]
  % three snapshots of 10 cells
  \foreach \r/\ttl in {0/{(a) Table $T_0=\mathrm{Ins}(a_1,\dots,a_4)$, before the push},-2.1/{(b) After \textbf{push} of $x$ (home 2; cells 2, 3 occupied; stored in cell 4): only cell 4 is written},-4.2/{(c) After \textbf{pop} of $x$: cell 4 is cleared, and the table is $T_0$ again}} {
    \node[font=\small\bfseries,anchor=west] at (-0.2,\r+0.95) {\ttl};
    \foreach \c in {0,...,9} \node[cell,fill=cFree] at (\c*0.9,\r) {};
    \foreach \c in {0,...,9} \node[lab,gray] at (\c*0.9,\r-0.55) {\c};
  }
  % occupied cells of T0
  \foreach \c/\k in {1/a_1,2/a_3,3/a_2,6/a_4} {
    \foreach \r in {0,-2.1,-4.2} \node[cell,fill=cIdx] at (\c*0.9,\r) {$\k$};
  }
  % x placed in cell 4 in row (b)
  \node[cell,fill=cUndo,draw=red!70!black,line width=1.2pt] at (4*0.9,-2.1) {$x$};
  \node[font=\small,anchor=west,text width=5.4cm,align=left] at (9.7,-2.1) {Insertion depends only on which cells are occupied and changes only one cell, so removing the \emph{newest} key restores the previous table bit for bit.};
  \draw[decorate,decoration={brace,amplitude=4pt,mirror}] (-0.45,-4.65) -- (8.55,-4.65) node[midway,below=5pt,font=\small]{identical in every cell: $\mathrm{Occ}(T)=\mathrm{Occ}(\mathrm{Ins}(\text{live sequence}))$ at all times};
\end{tikzpicture}

%%%%%%%%%%%%%%%%%%%%%%%%%%%%%%%%%%%%%%%% D3: dependency and orphans
\begin{tikzpicture}[x=1cm,y=1cm]
  \foreach \c in {0,...,9} \node[cell,fill=cFree] at (\c*1.05,0) {};
  \foreach \c in {0,...,9} \node[lab,gray] at (\c*1.05,-0.6) {\c};
  \foreach \c/\k in {2/u,3/v,4/w,5/y} \node[cell,fill=cIdx] at (\c*1.05,0) {$\k$};
  % probe path of y: home 2 -> cell 5
  \draw[arr,blue!60!black] (2*1.05,0.55) .. controls (2.9*1.05,1.35) and (4.1*1.05,1.35) .. (5*1.05-0.1,0.55);
  \node[lab,blue!60!black,above] at (3.5*1.05,1.18) {probe path of $y$: home 2, passes cells 2, 3, 4, stored in 5};
  \node[font=\small,anchor=north west,align=left,text width=10.2cm] at (-0.5,-1.0) {\textbf{Dependency:} $y\lhd x$ if the cell of $x$ lies strictly before the cell of $y$ on the probe path of $y$.\\Here $y\lhd u$, $y\lhd v$, $y\lhd w$. A deletion order is safe iff every key is cleared before the keys it depends on (linear extension of $\lhd$).};
  % unsafe: clear w first
  \begin{scope}[yshift=-5.3cm]
    \foreach \c in {0,...,9} \node[cell,fill=cFree] at (\c*1.05,0) {};
    \foreach \c/\k in {2/u,3/v,5/y} \node[cell,fill=cIdx] at (\c*1.05,0) {$\k$};
    \node[cell,fill=cBad,dashed] at (4*1.05,0) {};
    \node[lab,red!70!black,below] at (4*1.05,-0.45) {hole};
    \draw[arr,red!70!black] (2*1.05,0.55) .. controls (2.9*1.05,1.1) and (3.7*1.05,1.1) .. (4*1.05,0.55);
    \node[lab,red!70!black,above] at (3.4*1.05,1.0) {lookup of $y$ stops at the hole: $y$ is orphaned};
    \node[font=\small\bfseries,anchor=west] at (-0.5,1.85) {Clearing $w$ before $y$ (not stack order):};
  \end{scope}
  \node[font=\small,anchor=north west,align=left,text width=10.2cm] at (-0.5,-6.4) {\textbf{Stack order is always safe:} the newest key has no dependents, because a key's path crosses only cells that were filled before it. It is also the only order safe for every hash assignment (safe-orders theorem, part b). The number of keys orphaned by clearing one key is a queue length minus one (orphan law).};
\end{tikzpicture}

%%%%%%%%%%%%%%%%%%%%%%%%%%%%%%%%%%%%%%%% D4: flow diagram
\begin{tikzpicture}[x=1cm,y=1cm,node distance=0.8cm]
  \tikzset{dec/.style={diamond,draw,thick,aspect=2.6,font=\scriptsize,align=center,inner sep=0pt,fill=white},
           io/.style={draw,thick,rounded corners=6pt,font=\footnotesize,align=center,fill=cAcc,inner sep=4pt},
           ok/.style={draw,thick,rounded corners=6pt,font=\footnotesize,align=center,fill=cOk,inner sep=4pt},
           no/.style={draw,thick,rounded corners=6pt,font=\footnotesize,align=center,fill=cBad,inner sep=4pt}}
  \node[io] (s) at (0,0) {\texttt{declare(name)} arrives};
  \node[tb,below=0.55cm of s] (h) {compute home $h=F_K(\text{name})\bmod m$};
  \node[dec,below=0.55cm of h,minimum width=4.4cm,minimum height=1.25cm] (d1) {empty cell within\\$L$ probes of $h$?};
  \node[dec,below=0.75cm of d1,minimum width=4.4cm,minimum height=1.25cm] (d2) {record + undo entry\\fit in the buffer?};
  \node[tb,below=0.7cm of d2] (w) {write record, set index cell,\\push undo entry if the name hides another};
  \node[ok,below=0.55cm of w] (okn) {\texttt{kOk}};
  \node[no,right=2.6cm of d1] (pl) {\texttt{kProbeLimit}\\table unchanged};
  \node[no,right=2.6cm of d2] (nm) {\texttt{kNoMemory}\\table unchanged};
  \draw[arr] (s)--(h); \draw[arr] (h)--(d1);
  \draw[arr] (d1)-- node[right,lab]{yes} (d2); \draw[arr] (d2)-- node[right,lab]{yes} (w); \draw[arr] (w)--(okn);
  \draw[arr] (d1)-- node[above,lab]{no} (pl); \draw[arr] (d2)-- node[above,lab]{no} (nm);
  % exit scope
  \node[io] (e) at (-6.2,0) {\texttt{exit\_scope()}};
  \node[tb,below=0.55cm of e,text width=3.6cm] (e1) {pop undo entries of this scope (last in, first out)};
  \node[tb,below=0.55cm of e1,text width=3.6cm] (e2) {new name: clear its cell;\\hidden name: restore its record};
  \node[ok,below=0.55cm of e2,text width=3.6cm] (e3) {table equals the table before\\the scope was entered\\(exact undo)};
  \draw[arr] (e)--(e1); \draw[arr] (e1)--(e2); \draw[arr] (e2)--(e3);
  \node[font=\scriptsize,align=center,text width=3.6cm] at (-6.2,-7.9) {A refused request is all-or-nothing: no half-written state is ever visible.};
\end{tikzpicture}

%%%%%%%%%%%%%%%%%%%%%%%%%%%%%%%%%%%%%%%% D5: blocks and the refusal law
\begin{tikzpicture}[x=1cm,y=1cm]
  \def\R{2.9}
  \draw[thick] (0,0) circle (\R);
  % k = 6 empty cells at angles; blocks between
  \foreach \a/\g in {90/3,150/2,205/5,270/1,320/4,30/2} {
    \fill[white,draw=red!70!black,thick] (\a:\R) circle (0.2);
  }
  \foreach \a in {90,150,205,270,320,30} \node[font=\scriptsize] at (\a:{\R+0.55}) {};
  \draw[line width=3.4pt,cRec!85] (96:\R) arc (96:144:\R);   % block of 3? occupied run
  \draw[line width=3.4pt,cRec!85] (156:\R) arc (156:199:\R);
  \draw[line width=3.4pt,cRec!85] (211:\R) arc (211:264:\R);
  \draw[line width=3.4pt,cRec!85] (276:\R) arc (276:314:\R);
  \draw[line width=3.4pt,cRec!85] (326:\R) arc (326:384:\R);
  \draw[line width=3.4pt,cRec!85] (36:\R) arc (36:84:\R);
  \node[font=\small,align=center] at (0,0.25) {ring of $m$ cells};
  \node[font=\scriptsize,red!70!black,align=center] at (0,-1.0) {circles: the $k$ empty cells};
  \node[font=\footnotesize,align=center] at (0,-0.35) {cut at the $k=m-n$ empty cells};
  \node[font=\scriptsize,cRec] at (120:{\R+0.7}) {$g_1$};\node[font=\scriptsize,cRec] at (178:{\R+0.7}) {$g_2$};
  \node[font=\scriptsize,cRec] at (238:{\R+0.7}) {$g_3$};\node[font=\scriptsize,cRec] at (295:{\R+0.7}) {$g_4$};
  \node[font=\scriptsize,cRec] at (352:{\R+0.7}) {$g_5$};\node[font=\scriptsize,cRec] at (60:{\R+0.7}) {$g_6$};
  % legend/explanation
  \node[font=\small,anchor=west,align=left,text width=8.2cm] at (4.6,1.9) {\textbf{Block} = an empty cell and the run of occupied cells after it, of size $g_i$.\\[3pt]Blocks are i.i.d.\ Borel distributed (cycle lemma), conditioned on $\sum_i g_i=m$: weight $B(g)=g^{g-1}/g!$.};
  \node[font=\small,anchor=west,align=left,text width=8.2cm] at (4.6,-0.3) {A new key is refused iff its home starts a window of $L$ occupied cells. In a block of $g$ cells exactly $(g-L)^+$ homes do, so\\[2pt]\centerline{$R_m(n,L)=\frac{k}{m}\sum_{g>L}(g-L)\,\pi_g$.}};
  \begin{scope}[shift={(4.6,-2.6)}]
    \foreach \c in {0,...,9} \node[cell,minimum size=6mm,fill=cRec!85] at (\c*0.62,0) {};
    \node[cell,minimum size=6mm,fill=cFree] at (-0.62,0) {};
    \draw[thick,red!70!black] (-0.25,0.45) -- (-0.25,0.2); \node[lab,red!70!black,above] at (0.5,0.45) {bad homes: $g-L$ of them};
    \draw[decorate,decoration={brace,amplitude=3pt,mirror}] (-0.9,-0.45) -- (5.9,-0.45) node[midway,below,font=\scriptsize]{block of $g=11$ cells with $L=4$ (7 bad homes)};
  \end{scope}
\end{tikzpicture}

%%%%%%%%%%%%%%%%%%%%%%%%%%%%%%%%%%%%%%%% D6: persistent trunk with bounded churn
\begin{tikzpicture}[x=1cm,y=1cm]
  \tikzset{nd/.style={circle,draw,thick,minimum size=5mm,inner sep=0pt,fill=cIdx},nu/.style={circle,draw,thick,minimum size=4.2mm,inner sep=0pt,fill=cUndo}}
  \node[font=\small\bfseries,anchor=west] at (-0.4,4.7) {Scope tree: persistent trunk with bounded churn above it};
  \node[nd] (t0) at (0,0) {};
  \foreach \i/\n in {1/t1,2/t2,3/t3} \node[nd] (\n) at (\i*1.0,0) {};
  \draw[very thick] (t0)--(t1)--(t2)--(t3);
  \node[font=\scriptsize,below=3pt] at (1.5,-0.35) {trunk: $n$ names, never popped};
  \foreach \y/\lbl in {1.2/a,2.1/b,3.0/c} {
    \node[nu] (x\lbl1) at (4.0,\y) {}; \node[nu] (x\lbl2) at (5.0,\y+0.25) {};
    \draw[thick] (t3)--(x\lbl1)--(x\lbl2);
  }
  \node[nu] (xd1) at (4.0,-0.2) {}; \draw[thick] (t3)--(xd1);
  \node[font=\scriptsize,anchor=west,align=left] at (5.4,1.5) {fresh names pushed and popped\\at most $s$ live above the trunk\\(counting the one being pushed)};
  \draw[decorate,decoration={brace,amplitude=4pt}] (3.7,3.55) -- (5.4,3.55) node[midway,above=3pt,font=\scriptsize]{height $\le s$};
  % window panel
  \begin{scope}[shift={(0,-3.6)}]
    \node[font=\small\bfseries,anchor=west] at (-0.4,1.35) {Persistent part of the ring: a window of $L$ cells with at most $s-1$ empty cells};
    \foreach \c in {0,...,13} \node[cell,minimum size=6.4mm,fill=cRec!80] at (\c*0.7,0) {};
    \node[cell,minimum size=6.4mm,fill=cFree] at (4*0.7,0) {}; \node[cell,minimum size=6.4mm,fill=cFree] at (9*0.7,0) {};
    \draw[very thick,red!70!black] (-0.42,-0.42) rectangle (13*0.7+0.42,0.42);
    \node[font=\scriptsize,red!70!black,below] at (6.5*0.7,-0.5) {window of $L=14$ cells, $s-1=2$ empty cells};
    \node[font=\scriptsize,anchor=west,align=left,text width=6.6cm] at (10.2,0) {A chain of $s$ pushes fills the empty cells, then pushes at the window start: refused. Hence any refusal implies this window exists (churn bound).};
  \end{scope}
\end{tikzpicture}

%%%%%%%%%%%%%%%%%%%%%%%%%%%%%%%%%%%%%%%% D7: verification methodology
\begin{tikzpicture}[x=1cm,y=1cm]
  \tikzset{st/.style={tb,minimum width=3.0cm,minimum height=1.4cm,text width=2.8cm}}
  \node[st,fill=cAcc] (p) at (0,0) {\textbf{Statement and proof}\\pen and paper};
  \node[st] (e) at (3.7,0) {\textbf{Exhaustive enumeration}\\every home sequence, exact rationals};
  \node[st] (r) at (7.4,0) {\textbf{Real table}\\keyed hash, probe limit, z-test, lock-step};
  \node[st] (i) at (11.1,0) {\textbf{Independent derivation}\\Abel binomial route; separate referee pass};
  \node[st,fill=cOk] (o) at (5.55,-2.5) {\textbf{One script}\\\texttt{run\_all.sh}: stops on any mismatch};
  \node[st,fill=cOk] (f) at (9.25,-2.5) {\textbf{Tables and figures}\\written by the programs, none typed by hand};
  \draw[arr] (p)--(e); \draw[arr] (e)--(r); \draw[arr] (r)--(i);
  \draw[arr] (e.south) |- (o.west); \draw[arr] (r.south) -- (r |- o.north); \draw[arr] (o)--(f);
  \draw[arr] (i.south) |- (f.east);
\end{tikzpicture}
\end{document}
